% GPU Training Infrastructure Subsection for Chapter 4
% Created: November 2, 2025
% Documents computational infrastructure supporting BNN and deep learning models

\subsection{GPU Training Infrastructure for Neural Network Models}
\label{subsec:gpu-training-infrastructure}

Bayesian Neural Networks and deep ensemble methods require substantial computational resources that exceed typical CPU capabilities.
We developed a GPU training orchestration system to efficiently leverage NVIDIA RTX 4090 hardware (24GB VRAM, 16,384 CUDA cores) for parallel model training.

\paragraph{GPU Training Orchestrator.}
The orchestrator (\texttt{scripts/\allowbreak gpu\_training\_\allowbreak orchestrator.py}, 487 LOC) manages sequential training of compute-intensive models with automatic monitoring, checkpointing, and failure recovery.
Key features:
\begin{itemize}
  \item \textbf{Phase-based scheduling}: Critical models (Phase 1) run first, overnight batches (Phase 2) queued
  \item \textbf{Automatic checkpointing}: Saves model state every 10--50 epochs to resume from failures
  \item \textbf{Resource monitoring}: Tracks GPU memory, utilization, temperature via \texttt{nvidia-smi}
  \item \textbf{Git synchronization}: Auto-commits completed models to version control with metadata
  \item \textbf{Failure recovery}: Detects training failures, logs diagnostics, continues with next job
\end{itemize}

\paragraph{Training Job Definitions.}
\Cref{tab:gpu-training-jobs} summarizes the 8 GPU training jobs with estimated runtimes on RTX 4090:

\begin{table}[!ht]
\centering
\small
\caption[GPU training job schedule]{GPU training jobs on NVIDIA RTX 4090 (24GB VRAM, 16,384 CUDA cores) with estimated runtimes and phase priorities. Phase 1 jobs (critical models) execute first; Phase 2 (overnight batch) queued sequentially. Total runtime: 27.5 hours.}
\label{tab:gpu-training-jobs}
\begin{tabular}{llccr}
\toprule
\textbf{Job}  & \textbf{Model}  & \textbf{Phase}  & \textbf{Priority}  & \textbf{Est. Hours} \\
\midrule
1 & BNN Multi-Output & 1 & Critical & 2.5 \\
2 & GNN Ablation Study & 1 & Critical & 1.5 \\
3 & Attention BNN & 1 & Critical & 3.0 \\
4 & XGBoost GPU Tuning & 2 & Batch & 0.5 \\
5 & Massive BNN Ensemble & 2 & Batch & 8.0 \\
6 & cVAE Environment Model & 2 & Batch & 4.0 \\
7 & GNN Full Training & 2 & Batch & 2.0 \\
8 & Optuna Hyperparameter Search & 2 & Batch & 6.0 \\
\midrule
\textbf{Total (Phase 1)} & & & & \textbf{7.0} \\
\textbf{Total (Phase 2)} & & & & \textbf{20.5} \\
\textbf{Total (All)} & & & & \textbf{27.5} \\
\bottomrule
\end{tabular}
\end{table}

The orchestrator enables overnight batch training: queue all Phase 2 jobs before leaving the lab, return to 8 trained models with full logs and version control snapshots.

\paragraph{BNN Training Results (Phases 1--2).}
We completed BNN training through Phase 2 with Optuna hyperparameter optimization (\Cref{tab:bnn-training-results}):

\begin{table}[!ht]
\centering
\caption[BNN training results (Phases 1--2)]{Bayesian Neural Network training results from GPU orchestration system. Phase 1 baseline outperformed Phase 2 optimized model, suggesting overfitting in hyperparameter search or insufficient training data for complex architectures.}
\label{tab:bnn-training-results}
\begin{tabular}{lcc}
\toprule
\textbf{Metric}  & \textbf{Phase 1 (Baseline)}  & \textbf{Phase 2 (Optuna)} \\
\midrule
Architecture & 3-layer MLP & 3-layer MLP \\
Hidden Dims & [128, 64, 32] & [156, 89, 42] \\
Learning Rate & 0.001 & 0.0031 \\
Batch Size & 128 & 192 \\
Epochs & 200 & 200 \\
Training Time & 2.5 hours & 1.0 hour \\
Training Data & 2,163 games & 2,163 games \\
\midrule
MAE (Passing Yards) & \textbf{3.78} & 25.17 \\
TD Accuracy & \textbf{80.0\%} & 72.3\% \\
Model Size & 59 KB & 73 KB \\
Status & Success & Success \\
\bottomrule
\end{tabular}
\end{table}

\textbf{Analysis}: The Phase 1 baseline (3.78 yards MAE) significantly outperformed the Optuna-optimized model (25.17 yards), suggesting either (1) overfitting during hyperparameter search, (2) insufficient training data (2,163 games) for the optimized architecture's increased capacity, or (3) early stopping not properly tuned. This highlights the importance of careful validation splits and regularization in BNN training---automated hyperparameter search can degrade performance if not constrained by domain knowledge.

Future work should explore: larger training datasets (10k+ games), Bayesian optimization with priors on hyperparameter ranges, and cross-validated early stopping criteria.

\paragraph{Windows Platform Challenges.}
GPU training on Windows 11 encountered several platform-specific issues that required workarounds:

\begin{enumerate}
  \item \textbf{Unicode encoding crashes}: Windows default encoding (\texttt{cp1252}) caused crashes when logging UTF-8 team names (e.g., ``Buccaneers'' with accented characters in historical data). Fixed via:
  \begin{lstlisting}[language=Python,basicstyle=\ttfamily\scriptsize]
import sys
sys.stdout.reconfigure(encoding='utf-8')
sys.stderr.reconfigure(encoding='utf-8')
  \end{lstlisting}

  \item \textbf{Database schema mismatches}: Training scripts assumed \texttt{moneyline} column availability in the \texttt{games} table, but historical data (1999--2010) lacked moneylines (sportsbooks didn't systematically track them). Fixed by computing implied probabilities from spreads: $p \approx \Phi(\text{spread} / 13.5)$ where $\Phi$ is the standard normal CDF.

  \item \textbf{CUDA memory fragmentation}: RTX 4090's 24GB VRAM required batch size tuning to avoid OOM errors during large ensemble training. Implemented dynamic batch size reduction on CUDA OOM exceptions:
  \begin{lstlisting}[language=Python,basicstyle=\ttfamily\scriptsize]
try:
    model.fit(X_train, y_train, batch_size=batch_size)
except torch.cuda.OutOfMemoryError:
    batch_size //= 2  # Halve batch size, retry
    torch.cuda.empty_cache()
  \end{lstlisting}
\end{enumerate}

These Windows-specific issues consumed 2--3 hours of debugging time but are now documented in \texttt{docs/GPU\_TRAINING\_WINDOWS.md} for future researchers.

\paragraph{Integration with Ensemble.}
The BNN Multi-Output model integrates into the v3.0 ensemble (\Cref{sec:bayesian-hierarchical}, Ensemble Integration subsection) as the uncertainty quantification component. Posterior distributions from variational inference inform:
\begin{itemize}
  \item \textbf{Fractional Kelly sizing}: Lower posterior variance $\to$ higher stakes (model is confident)
  \item \textbf{Ensemble disagreement filtering}: High BNN uncertainty flags dubious predictions for manual review
  \item \textbf{CVaR constraint calibration}: BNN tail quantiles (5th, 95th percentiles) inform worst-case scenarios in risk-adjusted portfolio optimization
\end{itemize}

This GPU infrastructure reduced BNN training time from an estimated 12+ hours (CPU-only on MacBook M4) to 2.5 hours (RTX 4090), enabling rapid iteration on architectural variants and hyperparameter searches essential for player prop modeling. The orchestration system generalizes to any compute-intensive ML workflow requiring overnight batch processing with automatic recovery.

\paragraph{Future Enhancements.}
Planned improvements to the GPU training infrastructure include:
\begin{enumerate}
  \item \textbf{Multi-GPU support}: Parallelize independent models across multiple RTX 4090s for 4--8× speedup
  \item \textbf{Distributed training}: Use PyTorch Distributed Data Parallel (DDP) for single-model training across GPUs
  \item \textbf{Cloud integration}: Auto-provision AWS EC2 p4d.24xlarge instances (8× A100 80GB) for massive ensemble training
  \item \textbf{Hyperparameter search budget}: Implement intelligent search stopping based on marginal validation improvement
\end{enumerate}

The current single-GPU orchestrator serves as the foundation for these distributed extensions, with modular design enabling drop-in replacement of the task scheduler with Kubernetes or Slurm for HPC environments.
